% Codevorlagen
Silbentrennung auch bei Wörtern mit Bindestrich aktivieren:
SQL"=Funktionalitäten ("= wird zum Bindestrich)

Bilder einfügen:
% mögliche Positionen: h = here, p = extra page gemeinsam mit anderen Bildern, t = top, b = bottom, ! = zwingend
% Pfad relativ zu \graphicspath (main.tex), daher hier ohne "Ressourcen/Bilder/"-Präfix
\begin{figure}[hp]
    \centering
    \includegraphics[width=0.6\textwidth]{SolutionStructure.png}
    \caption{Struktur der Visual Studio Solution}\label{fig:VSSolution}
\end{figure}


Code einfügen (Caption/Label direkt am lstlisting, damit der Codeschnipsel im
Listingverzeichnis statt im Abbildungsverzeichnis erscheint):
\begin{lstlisting}[caption={Ausschnitt der abstrakten Klasse SQLQueryable}, label={lst:SQLQueryableAusschnitt}]
public abstract class SQLQueryable
{
		public abstract string TableName { get; }
		/// <summary>
		/// The column names of this table WITHOUT id
		/// </summary>
		protected abstract string[] ColumnNames { get; }
}
\end{lstlisting}

Caption ohne Zitat im Abbildungsverzeichnis:
\caption[Konfiguration des Ubuntu-Systems]{Konfiguration des Ubuntu-Systems, teilweise aus \cite{dockerUbuntuDockerDocs0200}}\label{fig:createMispUser}


tabularx-Tabelle (Spalte X füllt die Textbreite aus):
\begin{table}[hp]
    \centering
    \caption{Vergleich der Systemkomponenten}\label{tab:systemkomponenten}
    \begin{tabularx}{\textwidth}{l l X}
        \toprule
        Komponente & Version & Beschreibung \\
        \midrule
        Frontend   & 3.2.1   & Weboberfläche zur Visualisierung und Steuerung der Anwendung durch den Benutzer \\
        Backend    & 2.0.0   & REST-API-Server mit Geschäftslogik und Datenbankanbindung \\
        Datenbank  & 14.5    & Relationales Datenbankmanagementsystem zur persistenten Datenspeicherung \\
        \bottomrule
    \end{tabularx}
\end{table}


cleveref-Verweise (\cref{} für Kleinschreibung, \Cref{} für Großschreibung am Satzanfang):
Wie in \cref{fig:VSSolution} dargestellt, zeigt \cref{tab:systemkomponenten} die Übersicht.
\Cref{fig:VSSolution} zeigt die Struktur der Solution.
% Mehrere Referenzen auf einmal:
wie in \cref{fig:VSSolution,tab:systemkomponenten} zu sehen


Abkürzungen (glossaries-Package):
% Neue Abkürzung in Ressourcen/abkuerzungen.tex definieren:
% \newacronym{label}{Abkürzung}{Ausgeschriebene Bedeutung}
% Im Fließtext referenzieren, z. B.:
Ein \gls{ioc} zeigt eine mögliche Kompromittierung an.
% Erste Verwendung im Dokument schreibt automatisch "Indicator of Compromise (IOC)" aus,
% jede weitere Verwendung von \gls{ioc} zeigt nur noch "IOC".
% Groß-/Kleinschreibung:
\Gls{ioc} steht am Satzanfang groß.   % Satzanfang, erster Buchstabe groß
Man spricht hier allgemein von \GLS{ioc}s.   % komplett in Großbuchstaben, z. B. für Überschriften


Zitieren (biblatex, style=ieee):
% \parencite für eine abschließende/parenthetische Quellenangabe:
Dies wird bereits so beschrieben \parencite{amazonAmazonComputeSLA2022}.
% \textcite, wenn die Quelle das grammatikalische Subjekt des Satzes ist:
\textcite{amazonAmazonComputeSLA2022} beschreibt dieses Verhalten im Detail.
% \parencite rendert unter style=ieee als reines "[1]", \textcite dagegen als "Autor [1]"
% (schreibt den Autorennamen automatisch aus) — daher \textcite verwenden, wenn die Quelle als
% Satzsubjekt auftritt, \parencite für die abschließende/parenthetische Angabe. \footcite passt
% nicht zum nummerischen Klammerstil (Fußnote mit nur "[1]" wäre redundant) und wird deshalb hier
% nicht verwendet.


Deutsche Anführungszeichen (csquotes):
\enquote{Ein Beispieltext}
% \enquote{} setzt automatisch die korrekten Anführungszeichen für die aktive babel-Sprache
% (hier: ngerman, also „Ein Beispieltext"). Die manuelle Variante ``Ein Beispieltext'' erzeugt
% dagegen die falschen (englischen) Anführungszeichen und sollte nicht verwendet werden.


Links im Fließtext (hyperref/xurl):
\href{https://www.example.com}{Linktext}   % Link mit eigenem, lesbarem Anzeigetext
\url{https://www.example.com/ein/sehr/langer/pfad-der-sonst-nicht-umbrechen-wuerde}
% xurl (main.tex) erlaubt \url{}, an beliebigen Zeichen umzubrechen statt nur an "/" —
% wichtig bei langen URLs, damit sie nicht über den Seitenrand hinauslaufen.